\documentclass[11pt]{article}
\usepackage{geometry}
\geometry{a4paper, margin=1in}
\usepackage{amsmath, amssymb, amsthm}
\usepackage{mathtools}
\usepackage{enumitem}
\usepackage{hyperref}
\usepackage{tikz-cd}

\title{Cortical Neural Network}
\author{Alfred Ricker}
\date{February 2026}

\setlength{\parindent}{0pt}

\begin{document}

\maketitle
\section{Neurons}
Neurons are the fundamental unit of structure within the learning network. A single neuron is denoted as $h_i$.

\subsection{Activity}
The structure of the fundamental neuronal units is a simple ``activity ring'' $\mathcal{A}$. Each neuron $h_i$ has a time dependent value $\alpha(t) \in \mathcal{A}$ that is determined by the input connections and the \textit{kind} of connections that they are, as we shall see later.\\

Importantly, activity should contain a decay term, so that neurons don't remain in a static state until acted upon. If no messages are being sent to or from a neuron over a long period, the activity of this neuron should approach $0$, $\alpha(t) \to 0$ as $t \to \infty$.
\begin{equation}
    \alpha(t+1) = (1-\lambda)\alpha(t) + f
\end{equation}
Where $f$ is a term that is to be derived in later sections.

\subsection{Connection --- Incident and Terminal}
Neurons are connected to one another biologically through axons and dendrites. There is a directional flow of information. So if two neurons are connected, one is the transmitter and one is the receiver. To notate this, if $h_i$ transmits to $h_j$, we can write
\begin{equation}
    h_i \prec h_j
\end{equation}
We can say that $h_i$ is incident to $h_j$ and $h_j$ is terminal to $h_i$. The set of incident neurons of $h_i$ (neurons that transmit information to $h_i$) is to be denoted $P(h_i)$ (or just $P_i$) and the set of terminal neurons of $h_i$ is to be denoted $Q(h_i)$. Or, if we are dealing with a set of neurons $H$, we can take all the incident neurons of this set as $P(H) = \bigsqcup_{h \in H} P(h)$.\\

Connections serve the purpose of exciting or damping the activity of a given neuron. They can be thought of as weights. If $h_i \prec h_j$, we describe a connection between them by $w(h_i,h_j) \in \mathbb{R}$, where the incident neuron goes in the first argument and the terminal neuron goes in the second. The total synaptic input to neuron $h$ is
\begin{equation}
    f_h(t) = \sum_{p \in P(h)} \sigma(\alpha_p(t)) \cdot w(p, h)
\end{equation}
where $\sigma$ is the activation readout function $\sigma(x) = \frac{x}{|x|+1}$. This means the output signal of a neuron is also bounded in $(-1,1)$, consistent with the saturation philosophy. Combined with Eq.~\eqref{eq:update}, we have the full activation dynamics:
\begin{equation}\label{eq:full_dynamics}
    \alpha_h(t+1) = (1-\lambda)\alpha_h(t) + \frac{f_h(t)}{|f_h(t)|+1}, \qquad f_h(t) = \sum_{p \in P(h)} \sigma(\alpha_p) \cdot w(p,h)
\end{equation}

In some cases it is useful to split $P$ up into subsets $P = \bigcup_\mu P_\mu$ so that
\begin{equation}
    f_h(t) = \sum_\mu f_\mu(P_\mu)
\end{equation}
This will be important when different input streams (feedforward, recurrent, feedback) contribute to the same neuron through different mechanisms.

\subsection{Active Set and Sparsity}
A neuron $h_i$ is said to be \textit{active} at tick $t$ if its activation exceeds a threshold $\theta$:
\begin{equation}
    A(H, t) = \{h_i \in H : \alpha_i(t) > \theta \}
\end{equation}
We will impose sparsity constraints on certain neuron subsets, requiring that the active set remains small relative to the total population: $|A(H,t)| \approx k \ll |H|$.

\subsection{Representations as Population Patterns}
It is worth pausing to clarify a fundamental difference between this architecture and standard neural network models. In a transformer or convolutional network, each node carries an $n$-dimensional vector (an embedding, a hidden state, a feature map). The ``representation'' of a concept lives inside a single node as a high-dimensional object. Learning updates these vectors via backpropagation through the entire network.\\

In this formalism, each neuron carries a single scalar $\alpha_i(t) \in \mathbb{R}$. There are no embedding vectors. The representation of a concept is instead a \textit{pattern of activation across a population} --- the set $A(M, t) \subseteq M$ together with the activation values $\{\alpha_i(t) : m_i \in A(M, t)\}$. Two different concepts correspond to two different sparse subsets of the same neuron pool. The ``meaning'' is distributed across the population, not concentrated in any single node.\\

This has several consequences. First, the output of a region is not a vector but an activation pattern over the feed-out neurons $F_z$. If region $R_i$ needs to communicate to region $R_j$, the information is carried by the pattern $\{\sigma(\alpha_h(t)) : h \in F_z^{(i)}\}$, which the feed-in neurons $F_\omega^{(j)}$ receive through synaptic connections. If the output domain is external (e.g.\ motor commands), the map $z: R_Z \to Z$ reads out an activation pattern:
\begin{equation}
    z(\alpha) = \sum_{h \in F_z} \sigma(\alpha_h(t)) \cdot \phi_h
\end{equation}
where $\phi_h \in Z$ is a generator primitive associated with each output neuron (analogous to motor synergies or phonemes). This is a learned linear combination weighted by the current activation pattern.\\

Second, similarity between two representations is determined by the overlap of their active sets, not by a dot product of embedding vectors. If concept $A$ activates the set $a \subseteq M$ and concept $B$ activates $b \subseteq M$, their similarity is measured by $|a \cap b| / k$. The sparse coding constraint ($k \ll N$) ensures that randomly chosen patterns have near-zero overlap, while patterns learned from related inputs can share neurons.\\

Third, the capacity of the system scales with the combinatorial richness of sparse subsets rather than with the dimensionality of individual vectors. A region with $N = 1000$ neurons and sparsity $k = 20$ has $\binom{1000}{20}$ possible representations, which vastly exceeds what is achievable with a $1000$-dimensional real vector in practice.


\section{Neural Structures}
Before we get into the complex structure of regions $R$, it makes sense to cover some possible small scale neural structures, the possible ways to process information, and how smaller components may fit together to form regions.

\subsection{Translation Group}\label{subsec:rec}
Given a set of neurons $\{h_1, \ldots, h_n\}$ we can form a recurrent structure by letting
\begin{equation}
    h_1 \prec h_2 \prec \cdots \prec h_n \prec h_1
\end{equation}
This is a simple 1-dimensional case and can form the algebra of $\mathbb{Z}/n\mathbb{Z}$. A ring of $n$ neurons with cyclic connectivity supports a single ``active phase'' that advances under input, making it a natural substrate for modular position encoding. This leads to the question: how can a neural structure form a representation of an arbitrary dimensional translation group? One approach is to take a product of $L$ such rings with coprime periods $n_1, \ldots, n_L$, giving a representation of $\mathbb{Z}/n_1\mathbb{Z} \times \cdots \times \mathbb{Z}/n_L\mathbb{Z}$. By the Chinese Remainder Theorem, this encodes $\prod_\mu n_\mu$ unique positions using only $\sum_\mu n_\mu$ neurons.



\subsection{Input and Output}
To eventually form complex regional structures capable of representing a small ``brain-like'' container, we must have neurons that ``feed into'' the region through the intersections $\bigsqcup_\mu R_{j_\mu} \cap R_i$ and neurons that ``feed out'' into other regions $\bigcup_\nu R_{j_\nu} \cap R_i$. The ``feed in'' subset of $R_i$ can be defined as
\begin{equation}
    F_\omega = \{h \in R_i : \exists\, g \in P_h,\; g \notin R_i \}
\end{equation}
Which says that input neurons are those in which the set of incident neurons has elements outside of the region $R_i$. We can also define a set $B \subset F_\iota$ as the ``feedback'' neurons. These are those where $g \in P_h$ are also $g \in R_j$ where $R_j$ is interior to $R_i$. As for feed out neurons
\begin{equation}
    F_z = \{h \in R_i : \exists\, g \in Q_h,\; g \notin R_i \}
\end{equation}
These neurons function as transmission between regions, communicating the computed information from one region into the language of its neighbors. For adjacent regions $R_i, R_j$ where $R_j$ is deeper, the feedforward output of $R_i$ is embedded in the feedforward input of $R_j$: $F_\omega^{(i)} \subseteq F_\iota^{(j)}$.\\ 

Now we have identified a basic structure of communication within the network. What about external domains? A given domain that is capable of being processed by the network should have a defined map $\omega: \Omega \to \bigsqcup_i R_{\Omega_i}$. For a given region $R_{\Omega_i}$, there should be a surjection $\omega_i: \Omega_i \to R_{\Omega_i}$ where the activity pattern induced by some $x \in D$ is $\omega_i(x) = H \subset R_{\Omega_i}$.\\

Let's start with ways to define the input surjection, with the example of a 2-dimensional black and white image with pixel values $p \in [0,1]$ as is the case of an MNIST dataset. For a 28x28 image, this is a 2d array of size 784. We can initialize an input region for this pixel domain to have $n$ input neurons (the neurons at the boundary that receive the exterior data). Since $p \in  [0,1] \subset \Omega^{784}$, then we can have the simple map
\begin{equation}
    \alpha_{h_i} = p 
\end{equation}
We can have some overlap between regions, where region 1 covers pixels 1 through $n$, region 2 covers pixels $n-v$ to $2n -v$ and region $j$ covers $n-(j-1)v$ to $jn - (j-1)v$. Therefore the set of boundary neurons across several input regions communicates all information to the network. \\

As for an output map, let's start with assuming a classification task. Let the set of output neurons in an output region be denoted $F_z$ and the classification set be $Z = \{1,2,...,n\}$ (there is always a bijection between a set of $n$ classification objects and such a set $Z$). Therefore, we want 
\begin{equation}
    \bigsqcup_{F_\omega \subset R_{\mathrm{out}}} F_{\omega} \to Z 
\end{equation}

Let there set of output regions be $R_{z,\mu}$. The global classification is the consensus of these local populations. We define the global score $V(k)$ as the aggregate of local votes, typically a summation:
\begin{equation}
V(k) = \sum_{\mu} v_\mu(k) = \sum_{\mu} \sum_{h \in F_{z, \mu}} \sigma(\alpha_h(t)) \cdot w_{h, k}^\mu
\end{equation}
The system's prediction $\hat{z}$ is the class with the maximum consensus:
\begin{equation}
\hat{z} = \operatorname*{argmax}_{k \in Z} V(k)
\end{equation}



\subsection{$M$ Neurons}\label{subsec:m_neurons}
The model neurons $M = \{m_1, \ldots, m_N\} \subset R_i$ store object and feature representations as sparse activity patterns. At any tick $t$, only a small subset $A(M, t) \subset M$ should be active, with $|A(M,t)| \approx k \ll N$ and $k/N \approx 0.02$. This sparse distributed representation provides enormous combinatorial capacity: with $N = 1000$ and $k = 20$, there are $\binom{1000}{20} \approx 10^{41}$ distinguishable patterns.\\

An \textit{object model} $\mathcal{O}$ stored in a region is a set of (location, feature-pattern) pairs:
\begin{equation}
    \mathcal{O} = \{(w_\ell, a_\ell) : w_\ell \in \mathrm{States}(W),\;\; a_\ell \subseteq M,\;\; |a_\ell| = k \}
\end{equation}
Each pair encodes: ``when the where neurons indicate location $w_\ell$, the model neurons pattern $a_\ell$ should be active.'' The object model is not stored explicitly as a lookup table. It is implicit in the synaptic weight matrix $\{w(h_i, h_j)\}$ across the region. The collection of weights shapes an energy landscape whose attractors correspond to the sparse patterns $a_\ell$, conditioned on the location state $w_\ell$.

\subsubsection{M Neuron Activation}
The activation of a model neuron $m_i \in M$ receives three input streams: feedforward drive from $F_\omega$, a location gate from $W$, and recurrent drive from other $M$ neurons. These streams are combined through a gating mechanism. The total input is
\begin{equation}
    f_i^F(t) = \sum_{h_j \in F_\omega} w(h_j, m_i) \cdot \sigma(\alpha_j(t)), \qquad f_i^W(t) = \sum_{h_j \in W} w(h_j, m_i) \cdot \sigma(\alpha_j(t))
\end{equation}
\begin{equation}
    f_i^R(t) = \sum_{m_j \in M,\, j \ne i} w(m_j, m_i) \cdot \sigma(\alpha_j(t))
\end{equation}
The where neurons gate the feedforward input multiplicatively. The gating function $g: \mathbb{R} \to [0,1]$ is
\begin{equation}
    g(x) = \frac{x^2}{x^2 + \theta_W^2}
\end{equation}
so that the gate opens smoothly as the where signal $f_i^W$ grows. When the where signal is weak, $g \approx 0$ and the feedforward input is suppressed. When the where signal is strong, $g \approx 1$ and the feedforward input passes through. This implements a coincidence detection: only when the right sensory feature \textit{and} the right location signal are simultaneously present does the model neuron receive significant drive.\\

We also need lateral inhibition to enforce sparsity. Let $E_M(t) = \sum_{m_j \in M} \sigma(\alpha_j(t))$ be the total $M$ activity. Then the effective input to $m_i$ is
\begin{equation}\label{eq:m_input}
    f_i(t) = f_i^F(t) \cdot g(f_i^W(t)) + f_i^R(t) - \kappa \cdot \sigma(\alpha_i(t)) \cdot E_M(t)
\end{equation}
The inhibition term $\kappa \cdot \sigma(\alpha_i) \cdot E_M$ means that as the total population activity rises, it becomes harder for any individual neuron to grow, naturally enforcing a sparse active set and preventing runaway excitation. This $f_i(t)$ is then fed into the standard update rule Eq.~\eqref{eq:update}.

\subsubsection{Functional Subdivision}
Different subpopulations of $M$ may encode different feature types, so that $M = M_1 \cup M_2 \cup \cdots \cup M_p$ with sparse activation within each $M_q$: $|A(M_q, t)| \approx k_q$ with $\sum_q k_q = k$. The composite pattern $(A(M_1, t), \ldots, A(M_p, t))$ forms a richer representation than a single flat sparse code. Whether this subdivision should be architecturally imposed or allowed to emerge from learning is an open question.

\subsection{Where Neurons}
Where neurons, denoted $W \subset R_i$, describe ``position'' in an abstract sense (or concretely if modeling actual position within space). The where neurons should perform path integration similar to biological grid cells. We have the general condition that
\begin{equation}
    P(W) = \bigsqcup_j p_j \quad \text{where} \quad p_j \subset F_\omega
\end{equation}
Meaning each neuron $h \in W$ has an incident set of neurons in the feed in set $F_\omega$. Additionally, $W$ requires \textit{recurrent connectivity}: neurons in $W$ connect to other neurons in $W$. This internal recurrence is what maintains the location state across ticks and implements phase advance.
\begin{equation}
    \exists\, h_a, h_b \in W : h_a \prec h_b
\end{equation}

Most cortical columns will not be monitoring physical position directly. Instead, the ``where'' will often exist in some abstract sense. Therefore, the $W$ subset should be decomposed as
\begin{equation}
    W = W_\mathbb{T} \cup W_\mathbb{M}
\end{equation}
where $W_\mathbb{T}$ is a structurally hardwired grid component and $W_\mathbb{M}$ is a learned context component.

\subsubsection{The Grid Component $W_\mathbb{T}$}
The grid component consists of $L$ translation group modules (Section~\ref{subsec:rec}):
\begin{equation}
    W_\mathbb{T} = \bigsqcup_{\mu=1}^{L} W_\mu, \qquad W_\mu = \{h_{\mu,0},\; h_{\mu,1},\; \ldots,\; h_{\mu, n_\mu - 1}\}
\end{equation}
Each module $W_\mu$ has cyclic connectivity $h_{\mu,k} \prec h_{\mu, (k+1) \bmod n_\mu}$ and a current phase $\phi_\mu(t) \in \mathbb{Z}/n_\mu\mathbb{Z}$. At each tick, exactly one neuron per module carries high activation:
\begin{equation}
    \alpha_{\mu, k}(t) = \begin{cases} \alpha_{\max} & \text{if } k = \phi_\mu(t) \\ (1-\lambda)\,\alpha_{\mu,k}(t-1) & \text{otherwise} \end{cases}
\end{equation}
The phase updates by a displacement signal $\delta(t)$ received from a subset $D_W \subset F_\omega$:
\begin{equation}
    \phi_\mu(t+1) = \phi_\mu(t) + \delta_\mu \cdot \delta(t) \pmod{n_\mu}
\end{equation}
where $\delta_\mu \in \mathbb{Z}$ is the module's displacement sensitivity. If the module periods $\{n_1, \ldots, n_L\}$ are pairwise coprime, the full phase vector $(\phi_1, \ldots, \phi_L)$ uniquely encodes $\prod_\mu n_\mu$ positions by the Chinese Remainder Theorem. For example, $L = 3$ modules with periods $n \in \{5, 7, 11\}$ encode $385$ positions using only $23$ neurons.\\

The grid neuron activations are \textit{not learned} --- they are structurally determined by the ring connectivity and the displacement input. What \textit{is} learned are the weights from $W_\mathbb{T}$ to $M$, encoding which grid phases correspond to which features.

\subsubsection{The Context Component $W_\mathbb{M}$}
The context component $W_\mathbb{M}$ consists of neurons whose dynamics follow the standard activation update, with recurrent connections among themselves and input from $F_\omega$. These neurons learn to represent which ``space'' or reference frame the grid coordinates are operating in. The same grid phase can correspond to different absolute locations depending on the active context pattern.\\

For sensorimotor columns tracking physical position on an object, $W_\mathbb{T}$ does most of the work. For language or abstract reasoning columns where ``position'' lacks clean periodic structure, $W_\mathbb{M}$ becomes the primary location signal.


\section{Learning Rules}\label{sec:learning}
There are two classes of learning to consider: local Hebbian rules that operate on individual synapses based on pre- and post-synaptic activity, and region-wide modulatory signals (``enzymes'') that gate plasticity globally.

\subsection{Hebbian Weight Update}
For a connection $w(h_i, h_j)$ where $h_i \prec h_j$, the Hebbian update strengthens connections between co-active neurons:
\begin{equation}\label{eq:hebb}
    \Delta w(h_i, h_j) = \eta \cdot \sigma(\alpha_i(t)) \cdot \sigma(\alpha_j(t))
\end{equation}
The plasticity $\eta \in (0, 1)$ controls the learning rate. Since the $\sigma$ outputs are bounded in $(-1, 1)$, the maximum single-step update is bounded.\\

This rule applies to the learned connections: $F_\omega \to M$, $W \to M$, $M \to M$ (recurrent), and $M \to F_\omega$. The hardwired grid connections within $W_\mathbb{T}$ do not learn.

\subsection{Weight Decay}
To prevent saturation and allow the network to forget old associations, we include a slow weight decay:
\begin{equation}
    w(t+1) = (1 - \mu) \cdot w(t) + \Delta w(t), \qquad \mu \ll \eta
\end{equation}
The ratio $\mu / \eta$ controls the timescale separation between learning and forgetting. For stable long-term memory, we want $\mu \sim 10^{-4}$ while $\eta \sim 10^{-2}$.

\subsection{Modulatory Signal}
We may also have an ``enzyme'' induced learning rule, where some exogenous parameter $\nu$ modulates the plasticity. This could represent region-wide neuromodulatory signals (dopamine, norepinephrine, etc.) that globally increase or decrease the learning rate:
\begin{equation}
    \eta_{\text{eff}} = \eta \cdot \nu(t)
\end{equation}
where $\nu(t) \in [0, \nu_{\max}]$ is determined by signals from deeper regions or the global workspace. When a prediction error is large, $\nu$ increases, driving faster learning. When the region is in a familiar, well-predicted state, $\nu$ decreases toward baseline. The specifics of how $\nu$ is computed from network state are to be determined.





\section{Regions}
The cortical structure is to be thought of as a region $R_i$ that is given information transmission and processing structure. There are some preliminary definitions and axioms that I will use.
\begin{enumerate}
    \item An input region $R_\Omega$ is one with a (not necessarily injective or surjective) map from an external domain $\Omega$ to the region $\omega: \Omega \to R_\Omega$.
    \item An output region is one with the map $z: R_Z \to Z$.
    \item Two regions are connected if $R_i \cap R_j \ne \emptyset$.
    \item There are no disconnected regions. $\exists\, R_j : R_i \cap R_j \ne \emptyset \;\forall\, R_i \subset \mathcal{N}$.
    \item A path $\gamma$ is a sequence of distinct connected regions $R_1 R_2 \cdots R_n$. It is necessary that $R_i \cap R_{i\pm 1} \ne \emptyset$. However $R_i \cap R_j = \emptyset$ when $|j-i| > 1$ is allowed but not necessary.
    \item A region $R_i$ is interior to $R_j$ if the shortest path $\gamma_i$ between $R_i$ and some input region $R_\Omega$ is larger than the shortest input path $\gamma_j$. We can also say $R_j$ is exterior to $R_i$.
\end{enumerate}

To prevent the formation of isolated processing silos where regions only communicate in narrow chains ($R_{1} \to R_{1'} \to \dots$), we must impose a convergent hierarchy. This structure allows features from disparate sensory locations to be integrated into more complex, invariant representations.

\subsection{Hierarchy and Levels}
Let the set of all regions $\mathcal{R}$ be partitioned into discrete processing levels $\mathcal{L}_0, \mathcal{L}_1, \ldots, \mathcal{L}_D$.
\begin{enumerate}
    \item \textbf{Sensory Level ($\mathcal{L}_0$):} A region $R \in \mathcal{L}_0$ if and only if there exists a direct map from the external domain $\omega: \Omega \to R$.
    \item \textbf{Association Levels ($\mathcal{L}_k$):} For $k > 0$, regions in $\mathcal{L}_k$ receive their primary feedforward input from the feed-out neurons of regions in $\mathcal{L}_{k-1}$.
\end{enumerate}

\subsection{Convergence and Receptive Fields}
For any association region $R_j \in \mathcal{L}_k$, we define its \textit{receptive field} $\Gamma(R_j) \subset \mathcal{L}_{k-1}$ as the set of lower-level regions that project to it. To ensure information integration, we impose the \textbf{Convergence Axiom}:
\begin{equation}
    \forall R_j \in \mathcal{L}_k \quad (k>0), \quad |\Gamma(R_j)| \ge 2
\end{equation}
This ensures that no region simply ``passes along'' information from a single source; every association region must integrate features from multiple lower sources. The feed-in input to a neuron $h \in R_j$ is thus a summation over the feed-out sets of its receptive field:
\begin{equation}
    f_h^\text{feed}(t) = \sum_{R_i \in \Gamma(R_j)} \sum_{g \in F_z^{(i)}} w(g, h) \cdot \sigma(\alpha_g(t))
\end{equation}
This topology creates a ``pyramid of feature complexity.'' If $\Gamma(R_3) = \{R_1, R_2\}$, where $R_1$ detects vertical edges and $R_2$ detects horizontal edges, then $R_3$ is structurally capable of learning the representation for a corner, whereas $R_1$ and $R_2$ are not.

\subsection{Lateral Connectivity}
While the hierarchy integrates features vertically, regions must also resolve ambiguity horizontally. We define a lateral neighborhood $\mathcal{N}(R_i) \subset \mathcal{L}_k$ for regions within the same level.
\begin{equation}
    h \in F_\omega^{(i)} \implies P(h) \cap \left( \bigcup_{R_j \in \mathcal{N}(R_i)} F_z^{(j)} \right) \ne \emptyset
\end{equation}
Lateral connections are \textit{modulatory}. They do not drive a region to activation from rest, but rather lower the threshold for representations that are consistent with neighbors. This implements the distributed voting mechanism: if $R_i$ is ambiguous between objects $A$ and $B$, but neighbor $R_j$ strongly votes for $B$, the lateral input from $R_j$ will bias $R_i$ toward $B$.

\subsection{Top-Down Feedback}
The system must not be purely feedforward; higher-level context must inform lower-level processing (predictive coding). For every feedforward connection $R_i \prec R_j$ (where $R_i \in \mathcal{L}_{k-1}, R_j \in \mathcal{L}_k$), there exists a reciprocal feedback connection. \\

Feedback signals originate from the model neurons $M^{(j)}$ of the higher region and terminate on the feed-in neurons $F_\omega^{(i)}$ of the lower region. Unlike feedforward drive, feedback is \textit{depolarizing}:
\begin{equation}
    f_h^\text{total}(t) = f_h^\text{ext}(t) \cdot \left( 1 + \gamma \sum_{m \in M^{(j)}} w(m, h) \sigma(\alpha_m(t)) \right)
\end{equation}
where $\gamma$ is a gain factor. This multiplicative interaction means that when a lower-level feature is ``expected'' by the higher level, its gain is boosted, allowing it to win local inhibition competitions. Conversely, unexpected features are not suppressed, but simply lack this gain advantage. This allows the network to hallucinate expected features in noisy data while still remaining sensitive to strong, unexpected stimuli.


\end{document}